\footnotesize
\arrayrulecolor{black!22}
\begin{longtable}{L{0.055\linewidth} L{0.115\linewidth} L{0.195\linewidth} L{0.315\linewidth} L{0.19\linewidth} L{0.095\linewidth}}
\hline
\textbf{ID} & \textbf{Name} & \textbf{Metaphor} & \textbf{Mapping} & \textbf{Operational consequence} & \textbf{Boundary}\\
\hline
\endfirsthead
\hline
\textbf{ID} & \textbf{Name} & \textbf{Metaphor} & \textbf{Mapping} & \textbf{Operational consequence} & \textbf{Boundary}\\
\hline
\endhead
GM-TOOL-INTERFACE & Tool-interface principle & The CNNA tool locally exposes an active phase as a bright object against the complementary rest of the Lean codebase and its still unaudited semantic surroundings. & Bright sector = active phase/\allowbreak{}current object;\newline dark complement = remaining Lean code, open dependencies, unreviewed theory links, and unaudited semantics;\newline boundary/\allowbreak{}interface = Master-YAML, status, object register, proof dossiers, module/\allowbreak{}import DAG and later database views;\newline DtN/\allowbreak{}Schur-like reduction = extraction of relevant dependencies, modules, handoffs and review duties;\newline noise/\allowbreak{}uncertainty = semantic errors, missing blue markings, incomplete expert review or generated-artifact defects;\newline Birdsview = database/\allowbreak{}DAG/\allowbreak{}review view over the whole formal project, still secondary. & Each active phase must make its boundary to the remaining project explicit: dependencies, module interfaces, review duties, blue/\allowbreak{}reference markings and hidden complement risks must be visible before local closure is trusted. & This is an architecture metaphor for documentation and review practice, not a formal CNNA theorem, not an ontic claim, and not a replacement for Lean-primary verification or expert review.\\
\hline
\addlinespace[0.24em]
\end{longtable}
\arrayrulecolor{black}
\normalsize
